% GENERATED FILE. Edit canonical lessons, then run npm run generate:books.
% canonical-source-hash: fnv1a64:14161eee746b12a1
% canonical-lessons: KA-C469-kaicila, KA-C469-adalubadalu-madu, KA-C469-varanti, KA-C469-bekari, KA-C469-delivari

\chapter{Bags, Bakeries and Deliveries}
\label{ch:ka-bags-bakeries-and-deliveries}
\hlchaptermodality{469}

\begin{chapteropening}
\textbf{By the end of this chapter:} \emph{I can say a handbag, to swap, a warranty, a bakery and a home delivery.}

Complete the last lesson of chapter 469: I can say a handbag, to swap, a warranty, a bakery and a home delivery.
\end{chapteropening}

\section[\texorpdfstring{\kn{ಕೈಚೀಲ} (kaicīla)}{ಕೈಚೀಲ (kaicīla)}]{\kn{ಕೈಚೀಲ} (kaicīla) — a handbag}

\label{lesson:KA-C469-kaicila}



\begin{quote}
Before the new one: say the Kannada for a guarantee, then the Kannada for to fade.
\end{quote}

\subsection*{You'll want to know: \kn{ಕೈಚೀಲ}}
\textbf{\kn{ಕೈಚೀಲ}} — \emph{kaicīla} — \textquotedblleft{}a handbag\textquotedblright{}.

\begin{quote}
\kn{ಅಮ್ಮ} \kn{ತಮ್ಮ} \kn{ಕೈಚೀಲದಲ್ಲಿ} \kn{ಪರ್ಸ್} \kn{ಇಟ್ಟರು}.
\end{quote}

\begin{quote}
\emph{amma tamma kaicīladalli pars iṭṭaru.}
\end{quote}

\begin{quote}
Mother put her purse in her handbag.
\end{quote}

\kn{ಕೈ}, hand, plus \kn{ಚೀಲ}, a bag.

\begin{cousinweb}
\textbf{\kn{ಕೈಚೀಲ}} (\emph{kaicīla}) is a native Kannada compound of \emph{kai}, \textquotedblleft{}hand\textquotedblright{}, and \kn{ಚೀಲ}, \textquotedblleft{}a bag, a sack\textquotedblright{} (DEDR 2632).
\end{cousinweb}

\begin{grammarlens}[title={What you've built}]
One more word for getting around.
\end{grammarlens}

\subsection*{Your turn}
\begin{itemize}
\raggedright
  \item \emph{Say it:} \emph{kaicīla}
  \item \emph{Say it:} \emph{kaicīla}, once more
  \item \emph{Say it:} \emph{amma tamma kaicīladalli pars iṭṭaru.} — Mother put her purse in her handbag.
  \item \emph{Recall:} say the Kannada for a guarantee, then the Kannada for to fade, then say \emph{kaicīla} again
\end{itemize}

\begin{tcolorbox}[breakable,colback=teal!4,colframe=teal!35!black,title={Before you move on}]
What does \kn{ಕೈಚೀಲ} mean? (\textquotedblleft{}A handbag\textquotedblright{}.) Say it once more.
\end{tcolorbox}
\section[\texorpdfstring{\kn{ಅದಲುಬದಲು} \kn{ಮಾಡು} (adalubadalu māḍu)}{ಅದಲುಬದಲು ಮಾಡು (adalubadalu māḍu)}]{\kn{ಅದಲುಬದಲು} \kn{ಮಾಡು} (adalubadalu māḍu) — to swap; to exchange}

\label{lesson:KA-C469-adalubadalu-madu}



\begin{quote}
Before the new one: say the Kannada for to fade, then the Kannada for a handbag.
\end{quote}

\subsection*{You'll want to know: \kn{ಅದಲುಬದಲು} \kn{ಮಾಡು}}
\textbf{\kn{ಅದಲುಬದಲು} \kn{ಮಾಡು}} — \emph{adalubadalu māḍu} — \textquotedblleft{}to swap; to exchange\textquotedblright{}.

\begin{quote}
\kn{ನಾನು} \kn{ಚಿಕ್ಕ} \kn{ಅಂಗಿಯನ್ನು} \kn{ದೊಡ್ಡ} \kn{ಗಾತ್ರಕ್ಕೆ} \kn{ಅದಲುಬದಲು} \kn{ಮಾಡಿದೆ}.
\end{quote}

\begin{quote}
\emph{nānu cikka aṅgiyannu doḍḍa gātrakke adalubadalu māḍide.}
\end{quote}

\begin{quote}
I swapped the small shirt for a bigger size.
\end{quote}

Past (I …): \kn{ಅದಲುಬದಲು} \kn{ಮಾಡಿದೆ}, \emph{adalubadalu māḍide}.

An echo pair built on \kn{ಬದಲು}, instead.

\begin{cousinweb}
\textbf{\kn{ಅದಲುಬದಲು} \kn{ಮಾಡು}} (\emph{adalubadalu māḍu}) is built on \kn{ಬದಲು}, \textquotedblleft{}instead, in exchange\textquotedblright{}, usually traced to Arabic \emph{badal}, with a rhyming first half and the native verb \kn{ಮಾಡು}.
\end{cousinweb}

\begin{grammarlens}[title={What you've built}]
One more word for getting around.
\end{grammarlens}

\subsection*{Your turn}
\begin{itemize}
\raggedright
  \item \emph{Say it:} \emph{adalubadalu māḍu}
  \item \emph{Say it:} \emph{adalubadalu māḍu}, once more
  \item \emph{Say it:} \emph{nānu cikka aṅgiyannu doḍḍa gātrakke adalubadalu māḍide.} — I swapped the small shirt for a bigger size.
  \item \emph{Recall:} say the Kannada for to fade, then the Kannada for a handbag, then say \emph{adalubadalu māḍu} again
\end{itemize}

\begin{tcolorbox}[breakable,colback=teal!4,colframe=teal!35!black,title={Before you move on}]
What does \kn{ಅದಲುಬದಲು} \kn{ಮಾಡು} mean? (\textquotedblleft{}To swap; to exchange\textquotedblright{}.) Say it once more.
\end{tcolorbox}
\section[\texorpdfstring{\kn{ವಾರಂಟಿ} (vāraṇṭi)}{ವಾರಂಟಿ (vāraṇṭi)}]{\kn{ವಾರಂಟಿ} (vāraṇṭi) — a warranty}

\label{lesson:KA-C469-varanti}



\begin{quote}
Before the new one: say the Kannada for a handbag, then the Kannada for to swap.
\end{quote}

\subsection*{You'll want to know: \kn{ವಾರಂಟಿ}}
\textbf{\kn{ವಾರಂಟಿ}} — \emph{vāraṇṭi} — \textquotedblleft{}a warranty\textquotedblright{}.

\begin{quote}
\kn{ಹೊಸ} \kn{ಫೋನಿಗೆ} \kn{ಒಂದು} \kn{ವರ್ಷದ} \kn{ವಾರಂಟಿ} \kn{ಇದೆ}.
\end{quote}

\begin{quote}
\emph{hosa phōnige ondu varṣada vāraṇṭi ide.}
\end{quote}

\begin{quote}
The new phone has a one-year warranty.
\end{quote}

Shops say \kn{ವಾರಂಟಿ} for goods; \kn{ಖಾತರಿ} is the wider word.

\begin{cousinweb}
\textbf{\kn{ವಾರಂಟಿ}} (\emph{vāraṇṭi}) is a loan from English \emph{warranty}.
\end{cousinweb}

\begin{grammarlens}[title={What you've built}]
One more word for getting around.
\end{grammarlens}

\subsection*{Your turn}
\begin{itemize}
\raggedright
  \item \emph{Say it:} \emph{vāraṇṭi}
  \item \emph{Say it:} \emph{vāraṇṭi}, once more
  \item \emph{Say it:} \emph{hosa phōnige ondu varṣada vāraṇṭi ide.} — The new phone has a one-year warranty.
  \item \emph{Recall:} say the Kannada for a handbag, then the Kannada for to swap, then say \emph{vāraṇṭi} again
\end{itemize}

\begin{tcolorbox}[breakable,colback=teal!4,colframe=teal!35!black,title={Before you move on}]
What does \kn{ವಾರಂಟಿ} mean? (\textquotedblleft{}A warranty\textquotedblright{}.) Say it once more.
\end{tcolorbox}
\section[\texorpdfstring{\kn{ಬೇಕರಿ} (bēkari)}{ಬೇಕರಿ (bēkari)}]{\kn{ಬೇಕರಿ} (bēkari) — a bakery}

\label{lesson:KA-C469-bekari}



\begin{quote}
Before the new one: say the Kannada for to swap, then the Kannada for a warranty.
\end{quote}

\subsection*{You'll want to know: \kn{ಬೇಕರಿ}}
\textbf{\kn{ಬೇಕರಿ}} — \emph{bēkari} — \textquotedblleft{}a bakery\textquotedblright{}.

\begin{quote}
\kn{ಸಂಜೆ} \kn{ಬೇಕರಿಯಲ್ಲಿ} \kn{ಬಿಸಿ} \kn{ಬ್ರೆಡ್} \kn{ಸಿಗುತ್ತದೆ}.
\end{quote}

\begin{quote}
\emph{sañje bēkariyalli bisi breḍ siguttade.}
\end{quote}

\begin{quote}
In the evening you get hot bread at the bakery.
\end{quote}

Bread, biscuits, cakes and tea are sold there.

\begin{cousinweb}
\textbf{\kn{ಬೇಕರಿ}} (\emph{bēkari}) is a loan from English \emph{bakery}.
\end{cousinweb}

\begin{grammarlens}[title={What you've built}]
One more word for getting around.
\end{grammarlens}

\subsection*{Your turn}
\begin{itemize}
\raggedright
  \item \emph{Say it:} \emph{bēkari}
  \item \emph{Say it:} \emph{bēkari}, once more
  \item \emph{Say it:} \emph{sañje bēkariyalli bisi breḍ siguttade.} — In the evening you get hot bread at the bakery.
  \item \emph{Recall:} say the Kannada for to swap, then the Kannada for a warranty, then say \emph{bēkari} again
\end{itemize}

\begin{tcolorbox}[breakable,colback=teal!4,colframe=teal!35!black,title={Before you move on}]
What does \kn{ಬೇಕರಿ} mean? (\textquotedblleft{}A bakery\textquotedblright{}.) Say it once more.
\end{tcolorbox}
\section[\texorpdfstring{\kn{ಡೆಲಿವರಿ} (ḍelivari)}{ಡೆಲಿವರಿ (ḍelivari)}]{\kn{ಡೆಲಿವರಿ} (ḍelivari) — a home delivery; a delivery}

\label{lesson:KA-C469-delivari}



\begin{quote}
Before the new one: say the Kannada for a warranty, then the Kannada for a bakery.
\end{quote}

\subsection*{You'll want to know: \kn{ಡೆಲಿವರಿ}}
\textbf{\kn{ಡೆಲಿವರಿ}} — \emph{ḍelivari} — \textquotedblleft{}a home delivery; a delivery\textquotedblright{}.

\begin{quote}
\kn{ಅಂಗಡಿ} \kn{ಉಚಿತ} \kn{ಡೆಲಿವರಿ} \kn{ಕೊಡುತ್ತದೆ}.
\end{quote}

\begin{quote}
\emph{aṅgaḍi ucita ḍelivari koḍuttade.}
\end{quote}

\begin{quote}
The shop gives free home delivery.
\end{quote}

\kn{ಡೆಲಿವರಿ} \kn{ಹುಡುಗ}, a delivery boy; the verb is \kn{ತಲುಪಿಸು}, to deliver.

\begin{cousinweb}
\textbf{\kn{ಡೆಲಿವರಿ}} (\emph{ḍelivari}) is a loan from English \emph{delivery}.
\end{cousinweb}

\begin{grammarlens}[title={What you've built}]
One more word for getting around.
\end{grammarlens}

\subsection*{Your turn}
\begin{itemize}
\raggedright
  \item \emph{Say it:} \emph{ḍelivari}
  \item \emph{Say it:} \emph{ḍelivari}, once more
  \item \emph{Say it:} \emph{aṅgaḍi ucita ḍelivari koḍuttade.} — The shop gives free home delivery.
  \item \emph{Recall:} say the Kannada for a warranty, then the Kannada for a bakery, then say \emph{ḍelivari} again
\end{itemize}

\begin{tcolorbox}[breakable,colback=teal!4,colframe=teal!35!black,title={Before you move on}]
What does \kn{ಡೆಲಿವರಿ} mean? (\textquotedblleft{}A home delivery; a delivery\textquotedblright{}.) Say it once more.
\end{tcolorbox}
